\documentclass[a4paper,11pt]{article}

% ─── Geometry ────────────────────────────────────────────────────────────────
\usepackage{geometry}
\geometry{margin=1in, headheight=15pt}

% ─── Mathematics ─────────────────────────────────────────────────────────────
\usepackage{amsmath}
\usepackage{amssymb}
\usepackage{mathtools}

% ─── Graphics / Colours ──────────────────────────────────────────────────────
\usepackage{graphicx}
\usepackage{xcolor}
\definecolor{blue3}{RGB}{0,0,180}
\definecolor{green3}{RGB}{0,120,0}
\definecolor{orange3}{RGB}{200,100,0}
\definecolor{red3}{RGB}{180,0,0}
\definecolor{grey3}{RGB}{80,80,80}

% ─── Units ───────────────────────────────────────────────────────────────────
\usepackage{siunitx}
\sisetup{detect-all}

% ─── Links (hyperref BEFORE cleveref) ────────────────────────────────────────
\usepackage[colorlinks=true,linkcolor=blue3,citecolor=blue3,urlcolor=blue3]{hyperref}
\usepackage{cleveref}

% ─── Boxes ───────────────────────────────────────────────────────────────────
\usepackage{tcolorbox}
\tcbuselibrary{skins,breakable}

\newtcolorbox{pointsbox}[1][]{
    colback=grey3!8,
    colframe=grey3,
    fonttitle=\bfseries\small,
    #1
}

% ─── Tables / Lists ──────────────────────────────────────────────────────────
\usepackage{booktabs}
\usepackage{enumitem}

% ─── Notation commands ───────────────────────────────────────────────────────
\newcommand{\ktilde}{\tilde{k}}
\newcommand{\ytilde}{\tilde{y}}
\newcommand{\ctilde}{\tilde{c}}
\newcommand{\dd}{\mathrm{d}}
\newcommand{\bgp}{\mathrm{BGP}}

% ─── Header / Footer ─────────────────────────────────────────────────────────
\usepackage{fancyhdr}
\pagestyle{fancy}
\fancyhf{}
\lhead{\textbf{14E022 Macroeconomics I} \quad Practice Exam 4}
\rhead{BSE 2025--26}
\cfoot{\thepage}

% ═════════════════════════════════════════════════════════════════════════════
\begin{document}
% ═════════════════════════════════════════════════════════════════════════════

\begin{center}
    {\LARGE \textbf{14E022 Macroeconomics I}}\\[6pt]
    {\Large \textbf{Practice Exam 4}}\\[8pt]
    {\large BSE --- Academic Year 2025--26}
\end{center}

\vspace{8pt}
\noindent\textit{Instructions: the exam is made of 3 problems and a set of short questions totalling \textbf{100 points}.
You have \textbf{2 hours}. Show all derivations clearly. For short questions, concise and precise answers score highest.}
\vspace{4pt}

\noindent\rule{\textwidth}{0.4pt}

% ─────────────────────────────────────────────────────────────────────────────
\section{Question 1 --- Growth Accounting and Development Accounting \hfill (25 points)}
% ─────────────────────────────────────────────────────────────────────────────

Consider a Cobb-Douglas economy in which output of country $i$ is
\[
    Y_i = K_i^{\alpha}\,(h_i\,L_i)^{1-\alpha}\,A_i,
    \qquad \alpha = 1/3,
\]
where $K_i$ is physical capital, $L_i$ is the number of workers, $h_i$ is average human capital
per worker, and $A_i$ is TFP.  Two countries are observed:

\vspace{6pt}
\begin{center}
\begin{tabular}{lccc}
\toprule
Country & $Y_i/L_i$ (normalised) & $(K_i/Y_i)$ & $h_i$ \\
\midrule
US (benchmark) & 1.00 & 3.0 & 1.00 \\
Country D      & 0.20 & 2.4 & 0.60 \\
\bottomrule
\end{tabular}
\end{center}
\vspace{6pt}

\noindent All quantities are in units relative to the US benchmark.

\begin{enumerate}[label=\textbf{\arabic*.}, leftmargin=*, itemsep=8pt]

    \item \textbf{(8 points)} Derive the \emph{development accounting} decomposition that
    expresses output per worker $y_i = Y_i/L_i$ as a product of a capital-intensity term,
    a human-capital term, and TFP.  Using the data in the table, compute the fraction of the
    income gap between Country~D and the US that is explained by each factor.
    State clearly what fraction of the gap remains attributable to TFP.

    \item \textbf{(9 points)} Introduce firm-level heterogeneity following Hsieh and Klenow
    (2009).  Each firm $s$ faces an output wedge $\tau_{Y,s}$ and a capital wedge
    $\tau_{K,s}$.  Define revenue total factor productivity (TFPR) for a firm and explain
    why, under perfect competition with no distortions, TFPR would be equalised across all
    firms.  Show that dispersion in TFPR signals resource misallocation, and derive
    qualitatively how aggregate TFP falls below the frictionless benchmark when wedges are
    present.

    \item \textbf{(8 points)} ``Since TFP differences account for the bulk of cross-country
    income gaps, policies aimed at raising capital accumulation are largely irrelevant for
    development.''  Critically evaluate this claim.  Address whether high measured TFP is a
    cause or a consequence of higher income, identify one important identification challenge
    in interpreting TFP residuals causally, and discuss what the misallocation literature
    implies for development policy.

\end{enumerate}

\clearpage

% ─────────────────────────────────────────────────────────────────────────────
\section{Question 2 --- Baumol's Cost Disease and Structural Change \hfill (25 points)}
% ─────────────────────────────────────────────────────────────────────────────

Consider a two-sector economy with a \emph{progressive} sector~$P$ and a \emph{stagnant}
sector~$S$.  Output in each sector is produced using only labour:
\[
    Y_P = A_{P,t}\,L_P, \qquad Y_S = L_S,
\]
where productivity in the progressive sector grows at the constant rate
$\dot{A}_{P,t}/A_{P,t} = g > 0$, while productivity in the stagnant sector is fixed at~1.
A single competitive labour market equalises wages across sectors.  The total labour
endowment is $\bar{L}$ and is supplied inelastically.

Households have \textbf{Cobb-Douglas} preferences over the two goods:
\[
    U = Y_P^{\,\beta}\,Y_S^{\,1-\beta}, \qquad \beta \in (0,1).
\]

\begin{enumerate}[label=\textbf{\arabic*.}, leftmargin=*, itemsep=8pt]

    \item \textbf{(7 points)} Derive the equilibrium wage $w_t$ and the relative price
    $p_{S,t}/p_{P,t}$ as functions of $t$.  Show that, under Cobb-Douglas preferences, the
    labour allocation $(L_P/\bar{L},\, L_S/\bar{L})$ is constant over time.  Explain the
    economic mechanism that keeps expenditure shares constant even as relative prices rise.

    \item \textbf{(9 points)} Replace Cobb-Douglas preferences with non-homothetic
    preferences in which the income elasticity of stagnant services exceeds unity at high
    income levels.  Without solving a specific utility function, explain how the employment
    share of the stagnant sector evolves as income grows.  Carefully distinguish the
    \emph{demand channel} (rising income elasticity) from the \emph{supply channel}
    (Baumol cost disease), and discuss how the two channels interact to produce structural
    change.

    \item \textbf{(9 points)} Return to the Cobb-Douglas preference specification.  Derive
    the growth rate of aggregate output $Y = Y_P^{\,\beta}\,Y_S^{\,1-\beta}$.  Show
    formally that aggregate productivity growth lies strictly between 0 and $g$ regardless
    of the expenditure share $\beta$.  Explain why measured aggregate TFP growth
    systematically understates productivity growth in the progressive sector, and discuss the
    implication for interpreting cross-country productivity comparisons when countries differ
    in their sectoral composition.

\end{enumerate}

\clearpage

% ─────────────────────────────────────────────────────────────────────────────
\section{Question 3 --- Jones (1995) Semi-Endogenous Growth \hfill (25 points)}
% ─────────────────────────────────────────────────────────────────────────────

Consider a Romer-style economy with a continuum of differentiated intermediate goods of
measure $A_t$.  Final output is
\[
    Y_t = A_t^{\,\alpha/(1-\alpha)}\,L_{Y,t},
\]
where $L_{Y,t}$ is labour employed in final-goods production.  Knowledge (the measure of
blueprints) evolves according to
\[
    \dot{A}_t = \lambda\,A_t^{\,\phi}\,L_{A,t}, \qquad \phi \in (0,1),
\]
where $L_{A,t}$ is labour employed in R\&D and $\lambda > 0$ is research productivity.
Population grows exogenously at rate $n > 0$:
\[
    L_t = L_{Y,t} + L_{A,t}, \qquad \dot{L}_t/L_t = n.
\]

\begin{enumerate}[label=\textbf{\arabic*.}, leftmargin=*, itemsep=8pt]

    \item \textbf{(8 points)} On the balanced growth path, derive the long-run growth rate
    of $A_t$ (call it $g_A$) and the long-run growth rate of output per worker $Y_t/L_t$.
    Show explicitly that doubling the \emph{level} of population does not permanently raise
    $g_A$ --- the strong scale effect of Romer (1990) is absent.  Identify which
    parameter(s) determine long-run growth and explain the role of population growth~$n$.

    \item \textbf{(9 points)} Let $V_t$ denote the value of a new blueprint.  Write the
    free-entry condition for researchers and the no-arbitrage (Bellman) equation for $V_t$
    on the BGP.  Using the household Euler equation, derive the equilibrium share of labour
    devoted to R\&D, $\ell_A \equiv L_{A,t}/L_t$.  Show that $L_{A,t}$ grows
    proportionally to $L_t$ along the BGP --- the semi-endogenous growth property --- and
    verify consistency with your answer to question~1.

    \item \textbf{(8 points)} Compare Jones (1995) to Romer (1990).  In the Jones model, is
    long-run growth truly ``endogenous'' in any meaningful sense?  State the main
    theoretical objection to the Jones fix.  Briefly describe how fully endogenous
    alternatives (e.g.\ Peretto 1998 or Young 1998) eliminate scale effects without making
    growth depend on population growth, and discuss which set of predictions --- Jones or the
    fully endogenous alternatives --- appears more consistent with the available empirical
    evidence.

\end{enumerate}

\clearpage

% ─────────────────────────────────────────────────────────────────────────────
\section{Short Questions \hfill (25 points)}
% ─────────────────────────────────────────────────────────────────────────────

\noindent\textit{For these questions, short and sharp answers will receive the highest marks.
Aim for 3--6 sentences per question unless otherwise indicated.}

\vspace{8pt}

\begin{enumerate}[label=\textbf{\arabic*.}, leftmargin=*, itemsep=14pt]

    % ── SQ1 ──────────────────────────────────────────────────────────────────
    \item \textbf{(8 points) Dynamic Inefficiency and the Transversality Condition.}

    The Solow model can exhibit dynamic inefficiency --- capital accumulation beyond the
    Golden Rule level --- but the Ramsey model cannot.  Explain why the Ramsey model rules
    out dynamic inefficiency.  Define the Transversality Condition (TVC) precisely and show
    how it eliminates over-accumulation paths.  State the relationship between the TVC and
    the No-Ponzi condition on household borrowing.

    % ── SQ2 ──────────────────────────────────────────────────────────────────
    \item \textbf{(9 points) Level Effects, Growth Effects, and the Optimal Subsidy in
    Learning-by-Doing.}

    In a standard Ramsey economy, a permanent increase in the saving rate has only a
    \emph{level effect} on output per worker.  In the AK model, the same increase
    permanently raises the \emph{growth rate}.  Explain the mechanism behind this
    difference, with reference to the shape of the production function and the absence or
    presence of diminishing returns to capital.  In the Arrow (1962) learning-by-doing
    economy with aggregate production function $Y = AK$ (where $A$ embodies the
    externality), the decentralised growth rate is inefficiently low.  Derive the optimal
    subsidy to capital that restores the planner's allocation, expressing your answer in
    terms of the externality parameter $\alpha$ only, and interpret the result.

    % ── SQ3 ──────────────────────────────────────────────────────────────────
    \item \textbf{(8 points) Business-Stealing, Appropriability, and First-Best R\&D Policy.}

    In the Aghion--Howitt quality ladder model, equilibrium R\&D can be simultaneously too
    high \emph{and} too low relative to the social optimum.  Precisely define the
    business-stealing externality and the appropriability (consumer-surplus) externality.
    State the direction of each distortion, identify which tends to dominate in the basic
    model, and explain why implementing the first-best allocation requires more than one
    policy instrument.

\end{enumerate}

\vspace{20pt}
\noindent\rule{\textwidth}{0.4pt}

\begin{center}
    \textit{End of Practice Exam 4. Good luck!}
\end{center}

% ─────────────────────────────────────────────────────────────────────────────
\clearpage
\appendix
\section*{Appendix: Key Formulas Permitted in Exam}
% ─────────────────────────────────────────────────────────────────────────────

The following results may be used without re-derivation unless the question explicitly asks
you to derive them.

\subsection*{A.1 ~ Log-differentiation}
If $Z_t = X_t^a Y_t^b$, then $\gamma_Z = a\gamma_X + b\gamma_Y$.

\subsection*{A.2 ~ Ramsey / Euler Equation (standard)}
\[
    \frac{\dot{c}}{c} = \frac{1}{\sigma}(r - \rho - \sigma n),
    \qquad
    \frac{\dot{\ctilde}}{\ctilde} = \frac{1}{\sigma}(r - \delta - \rho - \sigma x)
\]
where $r = f'(\ktilde) - \delta$ is the net return to capital, $x$ is TFP growth, and
$n$ is population growth.

\subsection*{A.3 ~ Hamiltonian first-order conditions (current-value)}
Control $u$, state $x$, costate $\lambda$, discount $\rho$:
\[
    H_u = 0, \qquad \dot{\lambda} = \rho\lambda - H_x.
\]

\subsection*{A.4 ~ Cobb-Douglas factor payments}
$r K = \alpha Y$, $wL = (1-\alpha)Y$, so $r = \alpha Y/K$ and $w = (1-\alpha)Y/L$.

\subsection*{A.5 ~ Development accounting identity}
With $y_i = Y_i/L_i$, $\alpha = 1/3$, and capital-output ratio $\kappa_i = K_i/Y_i$:
\[
    \frac{y_i}{y_{US}} = \left(\frac{\kappa_i}{\kappa_{US}}\right)^{\alpha/(1-\alpha)}
    \cdot \frac{h_i}{h_{US}}
    \cdot \frac{A_i}{A_{US}}.
\]

\subsection*{A.6 ~ Jones (1995) BGP growth rate}
With $\dot{A}/A = \lambda A^{\phi-1} L_A$ and $L_A \propto L \propto e^{nt}$:
\[
    g_A^* = \frac{n}{1-\phi}.
\]

\subsection*{A.7 ~ CES monopoly markup}
With isoelastic demand elasticity $\varepsilon = 1/(1-\alpha)$:
\[
    p^* = \frac{1}{\alpha}, \qquad \pi = \frac{1-\alpha}{\alpha}\,x.
\]

\end{document}
